% P. Erdős, "Some asymptotic formulas for multiplicative functions", Bull. Amer. Math. Soc. 53 (1947), first page, reset in Mills 8A.
\documentclass[11pt]{article}
\usepackage{bulletin1947}
\begin{document}
\setcounter{page}{536}
\received{Received by the editors May 27, 1946, and, in revised form, December 11, 1946.}
\articletitle{SOME ASYMPTOTIC FORMULAS FOR\\ MULTIPLICATIVE FUNCTIONS}{p. erd\"os}

The present paper contains several asymptotic formulas for the sum of
multiplicative functions. A function $f(n)$ is called multiplicative if
$f(a\cdot b)=f(a)\cdot f(b)$ for $(a,b)=1$. We assume $f(n)>0$. In this paper
$f(n)$, $f_1(n)$ will always denote multiplicative functions. First we prove
the following theorem.

\par\smallskip\begingroup\textsc{Theorem 1.}\footnote{This result generalizes
a result of Wintner, Amer.\ J. Math.\ vol.~67 (1945) pp.~481--485.}\
\itshape Assume that the two series
\begin{equation}\tag{1}
-\sum_{p,\alpha}\frac{f(p^\alpha)-1}{p^\alpha}\,,\qquad
\sum_{p,\alpha}\frac{(f(p^\alpha)-1)^2}{p^\alpha},
\end{equation}
converge; then $f(n)$ has a mean value, that is,
\begin{equation*}
\lim_{n\to\infty}\frac{1}{n}\textstyle\sum f(m)
\end{equation*}
exists and is not equal to zero.\thmend

This result was conjectured in a slightly more special form at the end of my
paper \emph{Some remarks on additive and multiplicative
functions}.\footnote{Bull.\ Amer.\ Math.\ Soc.\ vol.~52 (1946) pp.~527--537.}

\textsc{Remark.} The convergence of (1) is the necessary and sufficient
condition for the existence of the distribution function of
$f(n)$.\footnote{P. Erd\"os and A. Wintner, Amer.\ J. Math.\ vol.~61 (1939)
pp.~713--721. See also footnote 2.}

For the sake of simplicity we assume $f(p^\alpha)=f(p)$. Then we prove
\begin{equation}\tag{2}
\lim_{n\to\infty}\frac{1}{n}\sum_{m\leqq n}f(m)=\prod_p\biggl(1+\frac{f(p)-1}{p}\biggr).
\end{equation}
It easily follows from (1) that the product on the right side of (2) converges
and thus the value of the limit is not 0.

We easily obtain from (1) that for every $\epsilon>0$
\begin{equation*}
\sum_{|f(p)-1|>\epsilon}\frac{1}{p}<\infty.
\end{equation*}
\end{document}
